\chapter{מנגנון תשומת הלב}

\section{המוטיבציה מאחורי תשומת הלב}

מנגנון תשומת הלב פתר בעיות קיימות בתרגום מכונה ובמשימות רצף־לרצף אחרות.
במודלים קודמים, כל וקטור הקלט היה דחוס לוקטור יחיד של "הקשר".
כתוצאה מכך, מודלים בעלי קיבולת קטנה היו מתקשים להתמודד עם קלטים ארוכים.

\section{הנוסחה הבסיסית}

תשומת הלב מחושבת על סמך שלוש מטריצות:
ערכי שאילתה \textenglish{Q}, ערכי מפתח \textenglish{K}, וערכי ערך \textenglish{V}.

הנוסחה הבסיסית היא:

\begin{equation}
\text{Attention}(Q, K, V) = \text{softmax}\left(\frac{QK^\top}{\sqrt{d_k}}\right)V
\end{equation}

כאן \(d_k\) היא המימד של וקטורי המפתח.

\section{שלבי החישוב}

חישוב תשומת הלב מתבצע בשלבים הבאים:

\begin{enumerate}
\item חישוב תאימות בין כל שאילתה לכל מפתח: \(QK^\top\)
\item קנה המידה על ידי חלוקה בשורש ממד המפתח: \(\frac{QK^\top}{\sqrt{d_k}}\)
\item החלת פונקציית \textenglish{softmax} להשגת משקלים: \(\text{softmax}(\cdot)\)
\item סיכום משוקלל של הערכים: \(\text{softmax}(\cdots) V\)
\end{enumerate}

\section{רעשי-דימיון עם מודלים קודמים}

למרות שתשומת הלב היא חדשה יחסית, ישנו דמיון לתורה חישובית של \textenglish{LSTMs}.
שתי הגישות מנסות ללמוד איזו אינפורמציה זורמת בין אסימונים.
ההבדל העיקרי הוא שתשומת הלב מאפשרת חישוב מקביל מלא, בעוד \textenglish{LSTM} דורשת עיבוד רציע.

\section{שרים מרובים}

בפרקטיקה, לא משתמשים בתשומת הלב יחידה אלא בכמה "ראשים" של תשומת הלב בו־זמנית.
כל ראש מלמד ליצור שאילתות, מפתחות וערכים שונים.
זה מאפשר למודל ללמוד התאמות מרובות וייצוגים עשירים יותר.

\section{יישום ו־אופטימיזציה}

בעת יישום \textenglish{Attention} בחומרה מודרנית,
החישוב משפר באמצעות טריקים כמו "הקבצה" של מטריצות גדולות
ושימוש בעיבוד מקביל \textenglish{GPU}.
זה מאפשר למודלים ענקיים להתרגל בזמן סביר.
